%!TEX root = main.tex
% ==============================================================================
% İSTANBUL ÜNİVERSİTESİ SOSYAL BİLİMLER ENSTİTÜSÜ
% BÖLÜM 1 (Madde 16/c, d ve Madde 17/b - İkinci Sistem)
% ==============================================================================
\chapter{KURAMSAL ÇERÇEVE VE KAVRAMSAL YAKLAŞIMLAR}

\section{Temel Kavramlar ve Yöntemsel İlkeler}

Sosyal bilim araştırmalarında kullanılan terimlerde ve üslupta birlik sağlanması yönergenin temel ilkelerindendir. Örneğin, metinde aynı anlama gelen terimlerin (örneğin ‘yasa’ / ‘kanun’ ya da ‘kuram’ / ‘teori’) hangisi seçildiyse o terimde tutarlılık gösterilmelidir.

Başka kaynaklara yapılan atıflar, araştırmacının tercihine göre metin içi yazar-yıl formatında \parencite[11]{sartori_demokrasi} veya sayfa altı referans dipnotu şeklinde verilebilir.\footnote{Giovanni Sartori, \textbf{Demokrasi Kuramı}, Çev. Deniz Baykal, Ankara, Siyasi İlimler Türk Derneği Yayınları, t.y., s. 11.} Yönerge Madde 19/a uyarınca, dipnot numaraları her bölüm başında yeniden 1'den başlar ve tek satır aralıklı, 10 punto ile dizilir.\footnote{Aynı esere yapılan sonraki atıflarda kısaltmalar (a.g.e., a.e. veya op. cit.) tutarlı olarak kullanılır: Sartori, \textbf{a.g.e.}, s. 22.}

\subsection{Doğrudan Alıntılar ve Blok Aktarma Kuralları}

Bir cümleyi aşmayan kısa alıntılar metin içinde çift tırnak (`"..."') ile gösterilir (Madde 18/a). Daha uzun aktarmalarda ise alıntı ile metin arasında belirgin bir ayrım yapılmalıdır (Madde 18/e ve Madde 4/b):

\begin{uzunalinti}
Kaynaklardan aynen yapılan alıntılar, bir cümleyi aşmayacak uzunluktaysa, tez metni içinde verilirler. Daha uzun alıntılarda, alıntının ilk ve son satırları ile tez metni arasında en az çift aralıklı boşluk bırakılarak, alıntı ile metnin ayırt edilmesi sağlanır. Bu tür alıntılar için, tez metninde kullanılan puntodan daha küçük punto kullanılır ve satırlar tek aralıklı yazılır (Madde 4/b, 18/e).
\end{uzunalinti}

\section{Özel Karakterler ve Transliterasyon}
Tezde yabancı alfabeler, hiyeroglif veya çivi yazısı transliterasyonları kullanılması durumunda uygun yazı tipleri devreye girer:
\begin{center}

\vspace{0.8em}

\translit{1-jr m ntw.f ḫr.tw sn.wy n wꜥ m(jw)t n wꜥ jt(j) jnpw rn pꜣ ꜥꜣ jw bꜣtꜣw rn pꜣ šrjw ḫr jr jnpw sw ẖry pr ẖry ḥmt {}
  2-jw pꜣy.f sn šrjw m-dj.f mj sḫrw n šrjw jw m ntw.f jrr.tw n.f ḥbsw jw.f ḥr šmyt m-sꜣ nꜣy.f jꜣwt r sḫtw \\{}
  3-[j]w m ntw.f jrr.tw skꜣ m ntw.f ꜥwꜣyw ḫr m ntw.f jrrt.n.f wpwt nb nty m sḫtw js.tw jw pꜣy.f sn šrjw} \\{}
\end{center}

\vspace{0.8em}

\begin{htranslit}
1-[LUGAL.GAL Ta-ba-ar]-⸢na⸣ m Ḫa-at-tu-ši-li LUGAL.GAL 
2-[LUGAL KUR URU Ḫa-at-]ti LÚ URU Ku-uš-šar KUR URUḪa-at-⸢ti⸣ 
3-[LUGAL e-ez-zi-a-at] ŠA fTa-wa-an-na-an-na DUMU ŠEŠ-ŠU  
4-[I-NA URU Š]a-na-wi-it-ta pa-it ša-an na-at-ta 
5-[ḫar-ni-ik]-⸢t⸣a na ut-ni-eš še-it ḫar-ni-ik-ta 
6-[nu-kán ERÍN.MEŠ] 2 AŠ-RA a-ša-an-du-la-ni da-<la>-⸢aḫ-ḫu⸣-un 
7-[ku-u]-e a-ša-u-wa-ar e-eš-[t]a 
8-[na-at A-⸢NA⸣] ERÍN.MEŠ a-ša-an-du-li pí-aḫ-ḫu-⸢un⸣ 
\tabletrule
9-[EGIR-an-da]-ma I-NA URU Za-la-pa pa-a[..............]-un |] \\
10-[na-an]⸢ḫa⸣r-ni-in-ku-un nu-uš-ši DMEŠ-ŠU ša-ra-a da-aḫ-ḫu-un|] \\
11-[Ù 3 GIŠ GIGIR.MEŠ MA-AD-NA-NU A-NA D.UTU URUTÚL-na pí-aḫ-ḫu-un|]
\tabletrule
12-[1 GUD KÙ.BABBAR 1 GEŠPU KÙ.BABBAR I-NA É D IM pí-aḫ-ḫu-un |] \\
13-[a-aš-ir-ma]-kán ku-i-e-eš na-aš I-NA É D Me-ez-zu-ul-la |] \\
14-[pí-da-aḫ-ḫu[...............................................]-un |]
  
\end{htranslit}

Bu sayede eskiçağ dilleri ve filolojik metinler özgün işaretleriyle bozulmadan aktarılabilmektedir.

